% ============================================================
%  Viyog: Separating Adversarial and Out-of-Distribution Inputs
%  via First-Layer Feature Norm Statistics
%  ESWEEK 2026 — IEEE Double-Column, Journal Track (TCAD)
%  IEEE double-column journal manuscript.
% ============================================================
\documentclass[10pt,letterpaper,journal]{IEEEtran}

\usepackage{cite}
\usepackage{amsmath,amssymb,amsfonts}
\usepackage{graphicx}
\usepackage{textcomp}
\usepackage[table]{xcolor}
\usepackage{booktabs}
\usepackage{subcaption}
\usepackage{stfloats}
\usepackage{url}
% Keep cross-reference/citation links, bookmarks, and metadata active while
% suppressing coloured boxes in the print version.
\usepackage[hidelinks]{hyperref}
\hypersetup{
  pdftitle={Viyog: Separating Adversarial and Out-of-Distribution},
  pdfauthor={A. Singh, A. Kumar, R. Kashyap, S. Park, A. Gurjar, Y. Ko, H. So, U. Hwang, K. Lee, A. Shrivastava},
  pdfkeywords={out-of-distribution detection, adversarial detection, embedded AI, first-layer feature statistics},
  bookmarksnumbered=true,
}
\usepackage{microtype}
\usepackage{array}
\usepackage{multirow}
\usepackage{algorithm}
\usepackage{algpseudocode}
\usepackage{makecell}
\usepackage{amsthm}
\usepackage{placeins}
\usepackage{pifont}
\newcommand{\cmark}{\ding{51}}% ✓
\newcommand{\xmark}{\ding{55}}% ✗

\newtheorem{proposition}{Proposition}

% ---- Float and caption separation for the 14-page format.
\setlength{\textfloatsep}{7pt plus 2pt minus 2pt}
\setlength{\floatsep}{7pt plus 2pt minus 2pt}
\setlength{\intextsep}{7pt plus 2pt minus 2pt}
\setlength{\dbltextfloatsep}{7pt plus 2pt minus 2pt}
\setlength{\dblfloatsep}{7pt plus 2pt minus 2pt}
\setlength{\abovecaptionskip}{3pt}
\setlength{\belowcaptionskip}{0pt}

% ---- Float placement parameters.
%      LaTeX's stock parameters are tuned for text-heavy papers with a couple
%      of figures. This paper has 18 figures + 6 tables, and the defaults
%      (textfraction=0.2 forcing >=20% text on every float page,
%      floatpagefraction=0.5 letting a float page sit half empty) were leaving
%      large gaps -- entire columns of whitespace around big figures. These
%      values let floats pack densely instead. Nothing about the content,
%      figures, or their order changes; only where LaTeX is permitted to put
%      them. floatpagefraction < topfraction is required to avoid placement
%      loops, and is respected here.
\renewcommand{\topfraction}{0.92}
\renewcommand{\bottomfraction}{0.85}
\renewcommand{\textfraction}{0.06}
\renewcommand{\floatpagefraction}{0.80}
\renewcommand{\dbltopfraction}{0.92}
\renewcommand{\dblfloatpagefraction}{0.80}
\setcounter{topnumber}{3}
\setcounter{bottomnumber}{2}
\setcounter{totalnumber}{4}
\setcounter{dbltopnumber}{3}

% ---- Compact reference-list spacing for the 14-page format.
\let\OLDthebibliography\thebibliography
\renewcommand\thebibliography[1]{%
  \OLDthebibliography{#1}%
  \setlength{\itemsep}{0pt}%
  \setlength{\parskip}{0pt}%
}

% ---- Legacy highlighting macro retained as a no-op for source compatibility. ----
\newcommand{\rev}[1]{#1}
\newcommand{\revtab}{}
\definecolor{viyoggreen}{HTML}{009E73}% Viyog brand green (matches figure palette)

% ---- Title --------------------------------------------------
\title{ Viyog: Separating Adversarial and Out-of-Distribution }
\author{Aman~Singh, Abhinav~Kumar$^{*}$, Rishab~Kashyap$^{\ddagger}$, Soyeong~Park$^{\ddagger}$, Ameya~Gurjar$^{\ddagger}$, Yohan~Ko, Hwisoo~So$^{*}$, Uiwon~Hwang, Kyoungwoo~Lee, and~Aviral~Shrivastava
\thanks{A. Singh, A. Kumar, R. Kashyap, A. Gurjar, and A. Shrivastava are with Arizona State University, Tempe, AZ, USA (e-mail: asing651@asu.edu; akuma294@asu.edu; rkashya6@asu.edu; agurjar2@asu.edu; aviral.shrivastava@asu.edu).}
\thanks{S. Park, Y. Ko, and K. Lee are with Yonsei University, Seoul, Republic of Korea (e-mail: psy980823@yonsei.ac.kr; Yohan.Ko@yonsei.ac.kr; kyoungwoo.lee@yonsei.ac.kr).}
\thanks{H. So is with Kyungpook National University, Daegu, Republic of Korea (e-mail: hwisoo.so@knu.ac.kr).}
\thanks{U. Hwang is with Ewha Womans University, Seoul, Republic of Korea (e-mail: uiwon.hwang@ewha.ac.kr).}
\thanks{$^{*}$Corresponding authors: Abhinav Kumar (e-mail: akuma294@asu.edu) and Hwisoo So (e-mail: hwisoo.so@knu.ac.kr).}
\thanks{$^{\ddagger}$These authors contributed equally to this work.}
}

\begin{document}
\maketitle

% ---- Abstract -----------------------------------------------
\begin{abstract}
\rev{Safety-critical systems face two different kinds of anomalous input, and
they call for opposite responses. \emph{Out-of-distribution} (OOD) inputs
should be abstained on. \emph{Adversarial} (ADV) attacks should be rejected
outright. Most anomaly detectors report only a single score, so they cannot
tell the two apart. We present \emph{Viyog}, a plug-in second stage that needs
no retraining and separates OOD from ADV using one simple signal: the
\emph{dormant-band roughness} of the first convolutional activation, i.e.\ how
spatially jagged its quietest channels are. Concretely, this is a
magnitude-normalised \emph{total variation} (the average change between
neighbouring pixels). It is read from the forward pass the model already
computes, and it stores a tiny $O(C)$ calibration profile, about 0.3\,KB for a
64-channel first layer, versus 4.5--40\,MB for the feature-distance baselines
we evaluate.

We find that, across the four $L_\infty$ attacks we test, adversarial inputs
leave more high-frequency residue than natural inputs do in these quiet
first-layer channels. Viyog exploits this asymmetry directly: it detects
adversarials at held-out AUROC \textbf{0.980} ($1.0$ is perfect, $0.5$ is
chance), compared with $0.659$ for the best confidence/logit score we
evaluate. Paired with a logit score, it forms a lightweight three-way
ID/OOD/ADV separator reaching balanced accuracy \textbf{0.798} on our
17-architecture held-out panel and \textbf{0.863} on GTSRB. We also prove that
the statistic is invariant to uniform activation rescaling in the
zero-offset limit, and give a conditional separation bound for affine layers.
Under an adaptive attack suite, a norm-targeted attack achieves \textbf{0\%
evasion}, though stronger TV-aware attacks do reduce detector AUROC. The
evaluation spans three image datasets, a 17-architecture held-out panel (three
of an original 20 architectures were excluded after a data-integrity check,
see Sec.~\ref{sec:setup}), four attacks, and a synthetic 1D sanity check.
Viyog adds no learned parameters and requires no extra forward pass; its
statistic costs a measured 2.28\% of CNN MACs and 3.9\% of inference
latency. Viyog is available on PyPI (\texttt{pip install viyog},
\href{https://pypi.org/project/viyog/}{pypi.org/project/viyog}) and as an
interactive live demo on Hugging Face
(\href{https://huggingface.co/spaces/amanyagami/viyog}{huggingface.co/spaces/amanyagami/viyog}).}
\end{abstract}

\begin{IEEEkeywords}
Out-of-distribution detection, adversarial robustness, feature norms,
deep neural networks, safety-critical systems, embedded AI, architecture
generalization.
\end{IEEEkeywords}

% =============================================================
\section{Introduction}
\label{sec:intro}

Deep neural networks (DNNs) now achieve state-of-the-art performance in
safety-critical domains: autonomous driving, medical imaging, embedded
vision. But a deployed system also needs to detect abnormal inputs and
respond to them appropriately before it acts on a prediction. Standard
classification pipelines do not do this~\cite{10.1145/3626314}.

Two kinds of anomalous input require \emph{opposite} responses
(Fig.~\ref{fig:problem}). \emph{OOD} inputs fall outside the training
distribution, so the right response is to abstain. \emph{Adversarial} (ADV)
inputs are crafted to force a misclassification, so the right response is to
reject them outright. Treating the two the same way is costly either way: it
wastes a defence on a benign input, or it lets an attack through silently.
Most lightweight baselines evaluated here do not draw this distinction at
all. We investigate a signal that does: one found in the network's
\emph{first} layer.

\begin{figure}[tb]\centering
\includegraphics[width=\columnwidth]{figs/rebuttal/fig_problem_concept.pdf}
\caption{\rev{\textbf{The core problem: OOD and ADV inputs require opposite responses, yet
both fool the classifier.}
\textbf{(1)~Input space:} an OOD (out-of-distribution) input lies outside the
training cluster; an adversarial (ADV) input is deliberately crafted to sit
\emph{inside} the cluster while forcing a wrong prediction.
\textbf{(2)~Classifier:} both can produce logits that the evaluated output-side
detectors do not reliably distinguish.
\textbf{(3)~Viyog:} reads the dormant-band roughness $V(x)$ of the first
convolutional activation map (a gradient-free statistic with $O(C)$ stored
state) and separates the two.
\textbf{(4)~Response:} OOD $\to$ abstain / flag for review; ADV $\to$ reject /
security alert.  Conflating them wastes adversarial defences on benign inputs or
silently accepts attacks.}}
\label{fig:problem}
\end{figure}

We identify an asymmetry in first-layer activation statistics that is consistent
across the evaluated models and can be read with low measured overhead.
For the evaluated digital attacks, the first convolutional or projection layer
contains a rough, broadband-like residue in quiet (``dormant'') channels, while
the tested ID and OOD inputs generally leave those channels smoother. Prior
work also reports early-layer and frequency-domain effects of adversarial
perturbations~\cite{FLOD,bakiskan2022early}. This is an empirical regularity of
our threat model, not a property guaranteed for every adversarial example.
\rev{An earlier, weaker form of Viyog exploited this asymmetry with the raw
first-layer $L_\infty$ norm (the \emph{$L_\infty$-form}): adversarials fall below
the ID mean and OOD at or above it. But the $L_\infty$-form is
\emph{regime-bounded}: strong only for \emph{far}-OOD (inputs far from the
training distribution, e.g.\ a different dataset) and near-chance for
\emph{near}-OOD (visually similar shifts) under a strict cross-model protocol.}
The primary signal is instead the \emph{relative roughness} of the dormant
channels: a magnitude-normalised total-variation (TV) statistic (Viyog,
Sec.~\ref{sec:method}) that measures \emph{shape} (spatial roughness), not
magnitude. \rev{In the zero-offset limit it is a $0$-homogeneous ratio, so
uniformly rescaling a feature map leaves it unchanged. This fact does not cover
general norm-preserving perturbations, which may change spatial roughness.} It is
complementary to logit OOD scores and yields three-way ID/OOD/ADV separation
(Sec.~\ref{sec:results}); the phenomenon persists across CNN and transformer
architectures and across $32\times32$ and $224\times224$ inputs. We ground this
in a formal Proposition (Section~\ref{sec:theory}) that gives a deterministic
bound for affine pre-activations under explicit smoothness, residue-roughness,
and residue-size conditions.

This paper makes five core contributions:

\noindent\textbf{C1: The Viyog scoring statistic.} A gradient-free
first-layer statistic (the magnitude-normalised total variation of the quietest
\emph{dormant} channels) that separates ADV from OOD where the raw $L_\infty$
norm cannot (mechanism confirmed on real ResNet-50 features, Fig.~\ref{fig:mechanism}). The continuous statistic $V(x)$ is itself the graded
degree-of-anomaly readout; the monotone routing squash is retained only to shape the binary operating curve, not as the severity scale. The raw
$L_\infty$ score is kept as the motivating baseline.

\noindent\textbf{C2: Theory for the shape statistic.} A formal
Proposition (Sec.~\ref{ssec:shape_theory}) proving that the dormant-band roughness
statistic is invariant to \emph{uniform rescaling} at zero stabilisation offset
and bounding its change under the exact affine residue
$e_c=w_c\!\ast\delta$. The separation result holds only when the stated
smoothness and rough-residue conditions hold for the affine pre-activation;
post-activation ReLU/GELU behavior is evaluated empirically rather than covered
by the bound. The ID/OOD/ADV ordering is therefore a conditional consequence,
not a worst-case guarantee; a general-affine operator-norm envelope
(Appendix~\ref{app:affine}) covers the raw-$L_\infty$ baseline.

\noindent\textbf{C3: Regime characterisation and a complementary
dormant-band detector.} Delimiting where the global $L_\infty$ mechanism
holds and recovering the lost near-OOD/strong-attack regime with Viyog, we
obtain a low-overhead three-way ID/OOD/ADV separation: across the
17-architecture held-out panel and two datasets Viyog is complementary to
logit OOD scores, and a lightweight supervised head over the full detector
panel reaches balanced accuracy \text{0.798}
(CIFAR-100, 17 archs) / \text{0.863} (GTSRB, 6 archs), the
logit$+$Viyog pair alone \text{0.758}/\text{0.768}, with $O(C)$ added state
(about 0.3\,KB at 64 channels), validated on ConvNeXtV2, ViT-B/16, and
Swin-T at $224\times224$.

\noindent\textbf{C4: Comprehensive empirical validation.} Experiments
across the dataset suite in Table~\ref{tab:coverage} and four $L_\infty$
attacks. Under the far-OOD
far-OOD single-direction protocol the raw $L_\infty$ score outperforms ten
established baselines by up to \textbf{$+$24.9~AUROC points}. The stricter
\emph{directionless} re-evaluation (which disqualifies any score whose high/low
meaning flips between models) over the 17-architecture held-out panel is
reported in Table~\ref{tab:main_results}.

\noindent\textbf{C5: Adaptive attack analysis.} In the evaluated attack
suite, the norm-targeted attack achieves $0\%$ evasion, a both-aware adaptive
attacker pays an $8\%$ success cost, and the stochastic
Expectation-over-Transformation (EOT) defence, which randomises which
dormant channels are read so the attacker cannot lock onto a fixed target,
restores the panel to \textbf{0.82} (Fig.~\ref{fig:adaptive_ladder}).

% =============================================================
\section{Background and Related Work}
\label{sec:related}

\subsection{OOD Detection Methods}
Table~\ref{tab:related} categorises representative prior work; surveys are
in~\cite{yang2024generalized,hong2024out}. Most evaluated detectors assign a
single anomaly score rather than an explicit OOD-versus-ADV decision. Szyc
et al.~\cite{szyc2023ood} show that rankings are protocol-sensitive, motivating
our multi-seed evaluation (Sec.~\ref{sec:setup}).

\begin{table}[tb]\centering\footnotesize
\caption{\rev{\textbf{Representative anomaly detectors and their OOD--ADV
separation capability in our evaluation.}
$\dagger$~Viyog scores both T2 and T3 above chance while remaining
gradient-free and state-minimal.}}
\label{tab:related}
\setlength{\tabcolsep}{3pt}
\resizebox{\columnwidth}{!}{%
\revtab
\begin{tabular}{@{}lllll@{}}
\toprule
Family & Examples & Score & Sep. OOD/ADV? & Cost \\
\midrule
Confidence & MSP, Energy, ODIN & logit & \xmark{} both high conf. & $O(1)$ \\
Distance & Maha, KNN, ViM & feat. dist. & \xmark{} ADV inside manifold & $O(D^2)$--$O(ND)$ \\
Uncertainty & MCD & post. var. & \xmark{} single score & $30{\times}$ fwd \\
Joint & UNICAD~\cite{unicad} & proto./DAE & partial & trained modules \\
\midrule
\textbf{Viyog}$^\dagger$ & \textbf{this work} & \textbf{shape $V$} & \textbf{\cmark{} T2\,0.980\,/\,T3\,0.844} & \textbf{$O(C)$} \\
\bottomrule
\end{tabular}}
\end{table}

\subsection{Adversarial Attacks, Defences, and Early-Layer Effects}
Adversarial attacks~\cite{szegedy2013intriguing} solve $\min_\delta\|\delta\|_p\;\text{s.t.}\;f(x{+}\delta,\theta)\neq y$.
Common attacks include FGSM~\cite{goodfellow2014explaining},
PGD~\cite{madry2017towards} (both $L_\infty$), and $L_2$-constrained
CW~\cite{carlini2017towards} and DeepFool~\cite{moosavi2016deepfool}.
A fundamental tension exists between robustness and accuracy~\cite{tsipras2018robustness,zhang2019theoretically}.

Defences against adversarial attacks fall into several families:
\emph{input preprocessing}~\cite{ziyadinov2023low},
\emph{feature denoising}~\cite{xie2019feature},
\emph{generative purification}~\cite{samangouei2018defense,nie2022diffusion,purify,silva2023diffdefense},
and \emph{adversarial training}~\cite{madry2017towards,zhang2019theoretically}.
Ensemble and Jacobian-based methods
provide additional robustness but at increased inference cost.
These defences do not generally target the OOD/ADV disambiguation problem; they
assume the threat class is known in advance, precisely the assumption
Viyog is designed to relax.

Early-layer effects are well documented: the first convolutional layer is
disproportionately critical for robustness~\cite{bakiskan2022early} and acts as
an implicit noise filter~\cite{FLOD} (the front-neuron concept formalised here);
pruning low-activation neurons improves robustness by forcing reliance on
high-magnitude responses~\cite{dhillon2018stochasticactivationpruningrobust},
and restoring suppressed activations~\cite{10640242} presupposes the suppression
Viyog detects. Filter-statistic~\cite{li2017adversarial} and
frequency-domain~\cite{wang2020high,wang2020towards} analyses likewise
distinguish adversarial perturbations from natural shift at early layers, but
none exploits the OOD/ADV asymmetry. These observations motivate the
conditional pre-activation proposition in Section~\ref{sec:theory}; the result
formalises sufficient roughness-separation conditions rather than universal
activation suppression.

\subsection{Joint OOD and Adversarial Detection}
A small body of work attempts the OOD--ADV separation directly, but at a cost
Viyog avoids.
Zhang et al.~\cite{zhang2025splitting} propose an OOD detector based on
adversarial gradient-attribution paths. It is not an OOD-versus-ADV router and
requires gradient computation at inference time. Feng et
al.~\cite{feng2022unified} propose a unified open-world detection perspective
but require joint training with anomaly labels.
Walkowiak et al.~\cite{walkowiak2023combining} combine outlierness scores and
feature extraction, while UNICAD~\cite{unicad} unifies attack detection and
novel-class identification in a single framework.
Detection methods based on feature statistics~\cite{li2017adversarial}
and disentangled representations have also been proposed.
Unlike gradient-attribution or
joint-training approaches, Viyog operates in a single gradient-free forward
pass and serves as a plug-in second stage atop any existing non-ID detector,
without architectural modification. The reported held-out protocol uses a
small labelled ID/OOD/ADV calibration split to fix score polarity and operating
points; the optional three-way head uses the same split
(Sec.~\ref{sec:results}). This efficiency
advantage is quantified in the system analysis that follows.

\subsection{System Efficiency and Deployment Motivation}
For embedded systems inference cost is a primary
constraint~\cite{10.1145/3626314}, yet prior detectors are expensive (ODIN and
gradient-norm backward passes, a full Mahalanobis covariance, MCD's 30 forward passes ($\sim\!30\times$ a single inference), diffusion denoising). Viyog adds zero parameters, no
gradient computation, and uses a first-layer activation hook plus a spatial
reduction (read as a dormant-band shape statistic);
the pipeline below motivates a candidate system integration.

In the Viyog two-stage pipeline (Fig.~\ref{fig:pipeline}), a standard inference
pass is followed by a first-stage non-ID detector, and Viyog attempts to
partition flagged inputs into OOD and ADV using the first-layer activation map
and the dormant-band roughness statistic $V$. This second stage provides a
candidate signal for choosing different responses; its errors and end-to-end
recall remain relevant in deployment (Sec.~\ref{sec:discussion}).

\begin{figure*}[tb]\centering
\begin{subfigure}[t]{0.44\textwidth}\centering
\includegraphics[width=\textwidth]{figs/rebuttal/fig_pipeline.pdf}
\caption{\rev{\textbf{Candidate Viyog two-stage pipeline.}
A normal forward pass drives a first-stage OOD gate (logit-based scores such as
Energy~\cite{NEURIPS2020_f5496252} or MSP (Maximum Softmax Probability~\cite{hendrycks2016baseline}))
and exposes the first-conv activation map $a=f_0(x)$ through a hook.
Flagged inputs enter \emph{Viyog}, which computes $V(x)$, the
dormant-band roughness (total-variation shape statistic) of $a$, and routes to
\textbf{ID} (pass), \textbf{OOD} (abstain / flag), or \textbf{ADV} (reject /
alert), adding $O(C)$ state (about 0.3\,KB at 64 channels), no extra forward
pass, and the measured
spatial-reduction cost reported in Sec.~\ref{sec:results}.}}
\label{fig:pipeline}
\end{subfigure}%
\hfill
\begin{subfigure}[t]{0.44\textwidth}\centering
\includegraphics[width=\textwidth]{figs/rebuttal/fig_mechanism.pdf}
\caption{\rev{\textbf{The Viyog mechanism confirmed on real ResNet-50 features
(CIFAR-100, $n{=}2000$ samples per class).}
\textbf{(a)}~Dormant-band channels (bottom 10\% alive, green) have the smallest
denominators, amplifying injected roughness.
\textbf{(b)}~ID and OOD co-cluster at low $V(x)$; ADV-PGD spikes high (AUROC
shown; dashed line~$=$ calibrated threshold).
\textbf{(c)}~Uniform-rescaling check (P1): doubling activation magnitude leaves
$V$ approximately unchanged at the chosen $\epsilon$; the adversarial example
has the elevated roughness ratio in this experiment.}}
\label{fig:mechanism}
\end{subfigure}
\end{figure*}

% =============================================================
\section{Theoretical Motivation and Analysis}
\label{sec:theory}

We formalise a sufficient condition under which a rough additive residue raises
the dormant-band statistic at an affine pre-activation
(Prop.~\ref{prop:shape}). The result proves uniform-rescaling invariance in the
zero-offset limit and conditional separation; it is not a worst-case guarantee
for arbitrary attacks or nonlinear post-activations. The argument is rooted
in the convolutional nature of the first layer shared by all architectures
evaluated: strided convolutions in ResNet~\cite{he2016deep},
DenseNet~\cite{huang2017densely}, and ConvNeXtV2~\cite{woo2023convnext};
and patch-embedding projections (which are mathematically equivalent to
strided convolutions) in ViT~\cite{dosovitskiy2020image} and
Swin-T~\cite{liu2021swin}. The insight that high-frequency and norm-altering
perturbations are concentrated at early layers is supported by prior
frequency-domain analyses~\cite{wang2020high,wang2020towards,ziyadinov2023low}.

\subsection{Why the First Layer?}
\label{ssec:rawbaseline}

\noindent\textbf{Natural-image smoothness (empirical premise A1).}
Natural image power spectra decay as $P(f)\!\propto\!f^{-2}$: low-frequency
components (slow gradients, textures) dominate; rapid pixel-to-pixel oscillations
are negligible.  Each first-layer channel $c$ produces a feature map
$a_c\in\mathbb{R}^{H\times W}$, the response of filter $w_c$ to input $x$.
For the tested natural inputs, $a_c$ is typically smooth: its
\emph{Total Variation} $\mathrm{TV}(a_c)$ (mean pixel-to-pixel jump, defined
formally in Eq.~\eqref{eq:vdef}) is small relative to the channel's mean
magnitude $\overline{|a_c|}$. Proposition~\ref{prop:shape} makes the required
per-channel inequality explicit; natural OOD inputs that violate it are a
failure mode rather than an exception to the premise.

\noindent\textbf{Rough adversarial residue (empirical premise A2).}
Gradient attacks add a perturbation $\delta$ to raise the loss. At an affine
first layer, each filter $w_c$ maps it to the exact additive residue
$e_c=w_c\!\ast\delta$. The four evaluated $L_\infty$ attacks usually produce
high relative TV in the selected channels. An $L_\infty$ constraint alone does
not impose this spectral property, however; a smooth or detector-aware
perturbation can violate A2.

\noindent\rev{\textbf{Why a normalised ratio, not a raw norm (0-homogeneity).}
A raw $L_\infty$ norm measures \emph{magnitude}, not \emph{shape}.  OOD images
can produce larger \emph{or} smaller first-layer activations than ID images, so
$\|a_c\|_\infty$ is direction-inconsistent; an attacker can further append a
norm-preserving penalty to keep $\|a_c'\|_\infty\!\approx\!\|a_c\|_\infty$ while
misclassifying.  The normalised ratio
$\widetilde{\mathrm{TV}}(a_c)=\mathrm{TV}(a_c)/(\overline{|a_c|}+\epsilon)$
(small $\epsilon>0$ for stability) addresses magnitude variation: TV and
$\overline{|\cdot|}$ are both 1-homogeneous, so their ratio is 0-homogeneous
when the stabilising offset is zero (P1 below). Thus uniform feature-map
rescaling does not change the idealised ratio. Preserving a norm while changing
spatial structure is not uniform rescaling and can change
$\widetilde{\mathrm{TV}}$; the adaptive experiments measure, rather than rule
out, that case.}




\subsection{The Dormant-Band Roughness Statistic}
\label{ssec:shape_theory}

\noindent\textbf{The statistic.}
Let $a = f_0(x) = W\!\ast x + b \in \mathbb{R}^{C\times H\times W}$ be the
first-layer activation map ($C$ channels, spatial size $H{\times}W$), with
$a_c \in \mathbb{R}^{H\times W}$ the feature map of channel $c$, filter $w_c$, and
bias $b_c$.  Define:
\begin{equation}
\begin{aligned}
V(x) &= \frac{1}{|\mathcal{B}|}\sum_{c\in\mathcal{B}}\widetilde{\mathrm{TV}}(a_c),
\qquad
\widetilde{\mathrm{TV}}(a_c)=\frac{\mathrm{TV}(a_c)}{\overline{|a_c|}+\epsilon},\\
\mathrm{TV}(a_c) &= \tfrac12\bigl(\overline{|\Delta_h a_c|}+\overline{|\Delta_w a_c|}\bigr),
\end{aligned}
\label{eq:vdef}
\end{equation}
where $\Delta_h a_c[i,j]{=}a_c[i{+}1,j]{-}a_c[i,j]$ and
$\Delta_w a_c[i,j]{=}a_c[i,j{+}1]{-}a_c[i,j]$ are pixel differences,
$\overline{u}=\tfrac1{HW}\sum_{i,j}u_{ij}$ is the spatial mean,
and $\epsilon=10^{-6}$ prevents division by zero.
The \emph{dormant band} is
$\mathcal{B}=\text{bottom-}p\%$ of \emph{alive} channels, where
$\mathcal{A}=\{c:\mathbb{E}_{\mathrm{ID}}\overline{|a_c|}>\kappa\}$
($\kappa=10^{-4}$) excludes permanently silent channels.

For the analysis below, $a$ denotes the affine pre-activation. A small
denominator can make the dormant band sensitive to roughness, but the
normalisation also makes dead-channel exclusion essential. Implementations that
hook a post-nonlinearity are evaluated empirically; they are not covered by the
affine proposition.
\rev{Fig.~\ref{fig:mechanism} confirms this on real ResNet-50 features: $V(x)$
separates by \emph{type of signal} (natural vs.\ attack), not by \emph{size} of
activation.}

\begin{proposition}[Conditional relative-roughness bound for an affine channel]
\label{prop:shape}
Fix a channel $c\in\mathcal B$ of the affine pre-activation and write
$a_c=w_c\!\ast x+b_c$, $\nu_c:=\overline{|a_c|}>0$. For
$x'=x+\delta$, affineness gives $a_c'=a_c+e_c$ exactly, where
$e_c=w_c\!\ast\delta$. Define
$r_c:=\overline{|e_c|}/\nu_c$,
$\rho_c:=\mathrm{TV}(e_c)/\overline{|e_c|}$ when
$\overline{|e_c|}>0$, and $\eta_c:=\epsilon/\nu_c$.
Assume:
\begin{itemize}\setlength\itemsep{1pt}
\item[\textbf{(A1)}] \emph{Natural smoothness:}
  $\mathrm{TV}(a_c)\le\beta_c\nu_c$; the same $\beta_c$ upper-bounds
  the natural ID/OOD comparison class.
\item[\textbf{(A2)}] \emph{Rough residue:} $\rho_c>0$ for the perturbation
  under consideration.
\item[\textbf{(A3)}] \emph{Nonzero residue:} $r_c>0$.
\end{itemize}
Then:
\begin{itemize}\setlength\itemsep{1pt}
\item[\textbf{(P1)}] \emph{Uniform-rescaling identity:} for $\lambda_c>0$,
\[
\widetilde{\mathrm{TV}}(\lambda_c a_c)
=\frac{\mathrm{TV}(a_c)}{\nu_c+\epsilon/\lambda_c}.
\]
Thus the statistic is exactly scale-invariant at $\epsilon=0$, and only
approximately so when $\epsilon/(\lambda_c\nu_c)$ is negligible.
\item[\textbf{(P2)}] \emph{Deterministic two-sided bound:}
\[
\frac{[\rho_c r_c-\beta_c]_+}{1+r_c+\eta_c}
\le \widetilde{\mathrm{TV}}(a_c') \le
\frac{\beta_c+\rho_c r_c}{|1-r_c|+\eta_c},
\]
where $[z]_+=\max(z,0)$. The upper bound becomes loose near $r_c=1$,
where cancellation can make the normalising denominator small.
\item[\textbf{(P3)}] \emph{Conditional separation:} if
\[
\rho_c r_c>\beta_c(2+r_c+\eta_c)
\]
for every $c\in\mathcal B$, then the perturbed pre-activation has larger
$V$ than any natural comparison satisfying A1. The per-channel margin is at
least
\[
\frac{\rho_c r_c-\beta_c(2+r_c+\eta_c)}
     {1+r_c+\eta_c}>0.
\]
\end{itemize}
P3 is a sufficient condition, not a guarantee that A1--A3 hold for an unseen
natural distribution or attack. ReLU is 1-Lipschitz, but that fact alone does
not preserve a bound on a normalised ratio; GELU is neither globally monotone
nor 1-Lipschitz. Consequently P2--P3 are not asserted for nonlinear
post-activation maps. Appendix~\ref{app:affine} instead bounds the local GELU
Taylor remainder, and Section~\ref{sec:results} reports the measured behavior.
\end{proposition}

\begin{proof}[Proof sketch (full proof in Appendix~\ref{app:affine})]
\textbf{P1.} Substitute the absolute 1-homogeneity of TV and
$\overline{|\cdot|}$ and divide numerator and denominator by $\lambda_c$.

\textbf{P2.} TV and the spatial $L_1$ norm satisfy triangle and reverse-triangle
inequalities:
\begin{align*}
\mathrm{TV}(a_c+e_c)&\le\mathrm{TV}(a_c)+\mathrm{TV}(e_c),\\
\mathrm{TV}(a_c+e_c)&\ge\mathrm{TV}(e_c)-\mathrm{TV}(a_c).
\end{align*}
Also,
$\nu_c|1-r_c|\le\overline{|a_c+e_c|}\le\nu_c(1+r_c)$.
Substituting $\mathrm{TV}(e_c)=\rho_c r_c\nu_c$, adding $\epsilon$ to
the denominator, and taking the nonnegative part of the lower numerator gives
P2.

\textbf{P3.} Any natural comparison satisfying A1 has normalised TV at most
$\beta_c$. Subtracting this from P2's lower bound gives the displayed margin;
averaging positive per-channel margins over the fixed band gives the result.
\end{proof}

\noindent\textbf{In plain words.} When a natural feature map is smooth and the
additive affine residue is sufficiently large and rough, the bound predicts a
higher $V$. The premise can fail for naturally rough OOD inputs or attacks that
optimise the spatial statistic, so the prediction must be checked empirically
under the intended threat model.

% =============================================================
\section{The Viyog Method}
\label{sec:method}

\subsection{Step-by-Step Derivation and Motivation}
The first design choice, reading at the first convolutional
layer, follows directly from the theory of Section~\ref{sec:theory}; the
remaining three (restricting to the dormant band, using a total-variation
\emph{shape} statistic, and held-out calibration) are developed in the
subsections that follow:

\noindent\textbf{First convolutional layer.} The first
convolutional layer has an exact affine perturbation relation before its
nonlinearity~\cite{bakiskan2022early,FLOD}, and, crucially, the adversarial residue
has not yet been transformed or pooled.  The shape statistic $V$ peaks at the
first conv (mean T2~$0.965$) and decays with depth in the six-model, two-attack
depth sweep; it is therefore the lowest-cost and highest-AUROC layer in that
experiment (Fig.~\ref{fig:shape_depth}).


Given this shape-based score, the routing rule is straightforward.

\subsection{Decision Rule and Downstream Routing}
Given a pool of non-ID inputs flagged by a first-stage detector, Viyog
applies a second-stage partition:
\begin{itemize}
  \item $\mathrm{Viyog}(x)>\tau$: \textbf{ADV}, apply adversarial
    defence, rejection, or security alert.
  \item $\mathrm{Viyog}(x)\le\tau$: \textbf{OOD}, apply abstention,
    fallback, or incremental-learning trigger.
\end{itemize}
A magnitude threshold allows precision--recall trade-offs. In the reported
protocol, score polarity and threshold $\tau$ are fixed from a labelled
ID/OOD/ADV calibration split and then applied once to disjoint test data
(Alg.~\ref{alg:viyog}); the end-to-end gate uses the same held-out protocol
(Sec.~\ref{ssec:e2e}).

\subsection{Dormant-Band Shape Statistic (Viyog)}
\label{ssec:viyogd}
\rev{The primary statistic $V(x)=\mathrm{Viyog}(x)$ uses the dormant band}
of Eq.~\eqref{eq:vdef}. Restricting it to \emph{active} channels is essential: up
to $25/64$ first-conv channels are permanently silent (e.g.\ DenseNet-121) and
would otherwise force a degenerate score. Across the 37-statistic battery $V$ tops
\emph{both} the ID-vs-ADV and OOD-vs-ADV rankings (Fig.~\ref{fig:sig_ranking}).

\begin{figure}[tb]\centering
\includegraphics[width=\columnwidth]{figs/rebuttal/fig_sig_ranking.pdf}
\caption{\rev{First-conv signature ranking by ID-vs-ADV (T2) AUROC on ResNet-50
(37 candidates), \textbf{coloured by family}. Magnitude/norm (orange), including
the raw $L_\infty$ baseline, sits near chance ($0.50$--$0.53$); shape (green,
TV/HF dormant-band) is best-in-class ($\approx0.98$) with \textbf{Viyog} first.
Distance (blue), logit (yellow), and distributional (purple) statistics fall
in between: shape, not magnitude, separates ADV from natural inputs here.}}
\label{fig:sig_ranking}
\end{figure}



\subsection{Algorithm and Complexity}
Fig.~\ref{fig:viyog_diagram} provides a step-by-step visual walkthrough
of the full Viyog pipeline with activation-map images at each stage.

\begin{figure*}[tb]\centering
\includegraphics[width=0.97\textwidth]{figs/rebuttal/fig_viyog_diagram.pdf}
\caption{\rev{\textbf{How Viyog works, in five steps (one forward pass, no extra
parameters).} \textbf{(1)}~Hook the first-conv activation map $a=f_0(x)$.
\textbf{(2)}~Rank channels by mean ID activation; the quietest $10\%$ form the
dormant band $\mathcal{B}$ (green). \textbf{(3)}~Read spatial roughness (total
variation) on $\mathcal{B}$: clean channels are smooth (low TV), adversarial
channels are jagged (high TV). \textbf{(4)}~Score $V(x)$, the mean normalized TV
over $\mathcal{B}$: the calibration split fixes polarity and threshold $\tau$.
\textbf{(5)}~Route the input: ID passes, OOD abstains, ADV is
rejected. Full method and routing rule in Sec.~\ref{sec:method}
(Alg.~\ref{alg:viyog}).}}
\label{fig:viyog_diagram}
\end{figure*}

\begin{algorithm}[tb]
\caption{Viyog: Dormant-band shape statistic for OOD-vs-ADV separation}
\label{alg:viyog}
\begin{algorithmic}[1]
  \Require Model $f$ with first affine layer $f_0$; labelled calibration sets
    $\mathcal{D}_{\rm ID},\mathcal{D}_{\rm OOD},\mathcal{D}_{\rm ADV}$;
    band fraction $\rho{=}0.1$; test sample $x$
  \State \textbf{Calibration} (once per model):
  \State \quad $p_c \leftarrow \frac{1}{|\mathcal{D}_{\rm ID}|}
    \sum_{x\in\mathcal{D}_{\rm ID}}
    \overline{|f_0(x)_c|}$ \Comment{per-channel ID mean activation}
  \State \quad $\mathcal{A} \leftarrow \{c : p_c > 10^{-4}\}$
    \Comment{drop permanently-dead filters}
  \State \quad $\mathcal{B} \leftarrow$ bottom-$\rho$ of $\mathcal{A}$ by $p_c$
    \Comment{dormant band}
  \State \quad $(d,\tau) \leftarrow$ fit polarity $d\in\{-1,+1\}$ and operating
    point on labelled calibration scores \Comment{then freeze}
  \State \textbf{Inference} (per sample, single forward pass):
  \State \quad $a \leftarrow f_0(x)$
    \Comment{hook; no extra pass, no gradients}
  \State \quad $s \leftarrow \frac{1}{|\mathcal{B}|}
    \sum_{c\in\mathcal{B}} \widetilde{\mathrm{TV}}(a_c)$
    \Comment{dorm-band roughness, Eq.~\eqref{eq:vdef}}
  \If{$d\,s > \tau$}
    \State \Return \textbf{ADV} (reject / defend)
      \Comment{high-frequency dorm-band residue}
  \Else
    \State \Return \textbf{OOD} (abstain / incremental learn)
  \EndIf
  \Statex \emph{Optional 3-way classifier:} fit a Linear Discriminant Analysis
    (LDA) on the labelled calibration split over the feature vector
    $[\text{logit score},\,V(x),\dots]$. This adds $<$0.1\,KB of state and
    returns a three-class ID/OOD/ADV decision (Sec.~\ref{sec:results}).
\end{algorithmic}
\end{algorithm}

Algorithm~\ref{alg:viyog} summarises the Viyog procedure: held-out calibration
computes the per-channel ID profile, dormant band $\mathcal{B}$, polarity, and
operating point; inference then adds only
$\mathcal{O}(|\mathcal{B}|\,h_1 w_1)$ operations and $O(C)$ stored state
($O(C)$; about 0.3\,KB at 64 channels), measured as 2.28\% of CNN MACs and 3.9\% of H200
inference latency, with four to five orders of magnitude less stored state than
the evaluated feature-distance detectors. We now validate it against eleven baselines across twenty
model--dataset configurations.

% =============================================================
\section{Experimental Setup}
\label{sec:setup}

\subsection{Datasets and Models}
All experiments use two primary ID datasets: \textbf{CIFAR-100}~\cite{krizhevsky2009learning}
($224{\times}224$, the main benchmark) and the safety-critical
\textbf{GTSRB}~\cite{stallkamp2012man} ($224{\times}224$).
Per-dataset OOD sets are drawn from SVHN~\cite{netzer2011reading}, iSUN, LSUN,
and Places365; MNIST serves as a supplementary near-OOD probe for DenseNet.

The model panel comprises \textbf{17 architectures} spanning every major
design family: ResNet-50/101/152~\cite{he2016deep},
DenseNet-121/161/169/201~\cite{huang2017densely},
ConvNeXtV2~\cite{woo2023convnext}, ViT-B/16~\cite{dosovitskiy2020image},
Swin-T~\cite{liu2021swin}, MobileNetV3-L, MobileNetV4-M,
EfficientNet-Lite0, EfficientNetV2-L,
EfficientViT-B1, FastViT-SA12, and EdgeNeXt-Small, all
fine-tuned from ImageNet pretraining in
PyTorch~\cite{paszke2019pytorch} without adversarial augmentation. (ResNet-18/34
and MobileOne-S1 were excluded: their archived first-layer features are
identically zero, Table~\ref{tab:limits}.)
Training-seed confidence intervals (${\le}3$ seeds per architecture) are
reported throughout; per-architecture std on T2 is ${\le}0.004$ (Fig.~\ref{fig:seed_forest}).

\subsection{Attacks, Baselines, and Metrics}
Four $L_\infty$ attacks span the threat model: FGSM~\cite{goodfellow2014explaining},
BIM~\cite{kurakin2016adversarial}, PGD~\cite{madry2017towards}, and
APGD-CE~\cite{croce2020reliable} via TorchAttacks~\cite{kim2020torchattacks} at
$\varepsilon{=}8/255$. Eleven \texttt{pytorch-ood}~\cite{kirchheim2022pytorch}
baselines, \textbf{MSP}, \textbf{Entropy}, \textbf{MaxLogit}, \textbf{Energy},
\textbf{GEN}, \textbf{ODIN} ($T{=}1000$), \textbf{MCD} (30 passes), \textbf{Mahalanobis},
\textbf{KNN}, \textbf{OpenMax}, and \textbf{KL-Matching}, are compared on
\textbf{AUROC}, \textbf{FPR@95TPR} (false-positive rate at 95\% true-positive rate), and \textbf{AUTC} (area under the threshold curve, critical for deployment where post-hoc tuning is infeasible).
We evaluate three tasks: \textbf{T1}~ID-vs-OOD, \textbf{T2}~ID-vs-ADV (the
security task), and \textbf{T3}~OOD-vs-ADV (the second-stage routing task).

\noindent\textbf{Held-out calibration protocol.} For the main CIFAR-100
evaluation, each pool is partitioned deterministically into 20\% calibration
and 80\% test data (split seed 0). Calibration data alone determine the
alive/dormant channels (absolute threshold $p_c>10^{-4}$), KL-Matching
references, score polarity, and operating points. Polarity and operating points
may use labelled ID/OOD/ADV calibration samples; they are then frozen and
evaluated once on the disjoint test split. Reported AUROCs use the
calibration-selected polarity, not a test-set oracle or
$\max(\mathrm{AUROC},1-\mathrm{AUROC})$. A $\times$ in
Table~\ref{tab:main_results} indicates that the raw score's high/low direction
varies across architectures; it is a diagnostic rather than a post-hoc test
orientation rule~\cite{szyc2023ood}.
Table~\ref{tab:coverage} studies marked $\dagger$ predate the exclusion above.

\rev{\noindent\textbf{Experimental coverage.}
Table~\ref{tab:coverage} lists the experiments. The main detection sweep uses
the held-out protocol above; specialised studies use the protocols stated in
their sections. Beyond the main detection sweep we add a
layer-depth sweep (Fig.~\ref{fig:shape_depth}), a \emph{measured} on-silicon cost
benchmark against all eleven baselines
(Figs.~\ref{fig:baselines_cost},~\ref{fig:onsilicon}), an adaptive
ladder with an EOT stochastic-band defence (Sec.~\ref{sec:adaptive},
Fig.~\ref{fig:adaptive_ladder}), training-seed CIs
(Fig.~\ref{fig:seed_forest}), and a 1D non-vision sanity check
(Appendix~\ref{app:nonvision}). The adaptive ladder and training-seed CIs were
run on the original 20-architecture panel \emph{before} the MobileOne-S1/
ResNet-18/ResNet-34 data-integrity issue was identified and have not been
rerun against the corrected 17-architecture panel for this submission
(Table~\ref{tab:limits}).}

\begin{table}[tb]\centering\footnotesize
\setlength{\tabcolsep}{3.5pt}
\renewcommand{\arraystretch}{1.1}
\caption{\rev{\textbf{Experimental coverage.} Main detection uses a
deterministic 20\% calibration / 80\% test split. $\dagger$ predates the
data-integrity fix (Table~\ref{tab:limits}), not rerun on the 17-arch panel.}}
\label{tab:coverage}
\revtab
\begin{tabular}{@{}p{2.45cm}p{3.15cm}p{1.85cm}@{}}
\toprule
Experiment & Coverage & Measures \\
\midrule
Main T1/T2/T3 & CIFAR-100, 17 usable archs, 20/80 held-out split & AUROC, FPR@95, recall \\
Measured cost & CIFAR-100, H200 (CUDA events) & latency, mem, energy \\
Layer-depth sweep & 6 archs $\times$ 2 attacks $\times$ all depths & T2/T3 vs.\ depth \\
Adaptive A0--A3, EOT\textsuperscript{$\dagger$} & 20 archs; norm/TV-aware $+$ stochastic & evasion \%, AUROC \\
Strong adaptive A4 & 6 archs (C100$+$GTSRB), held-out basis & detector floor \\
Training-seed CIs\textsuperscript{$\dagger$} & 20 archs $\times$ ${\le}3$ seeds & T2/T3 std \\
Non-vision sanity & 1D signals, FGSM/PGD & T2 AUROC \\
\bottomrule
\end{tabular}
\end{table}

% =============================================================
\section{Results and Discussion}
\label{sec:results}

\subsection{Detection Quality Across 17 Architectures}
\rev{With the evaluation protocol of Sec.~\ref{sec:setup} fixed, we first report
raw detection quality across all 17 usable backbones.}

\begin{figure}[tb]
  \centering
  \includegraphics[width=0.88\columnwidth]{figs/rebuttal/fig_detection_quality.pdf}
  \caption{\rev{Detection quality across the 17
    CIFAR-100 architectures, by task and method family. Viyog leads on both security tasks: \textbf{T2~0.980} (ID-vs-ADV) and \textbf{T3~0.844} (OOD-vs-ADV), while logit detectors
    win the T1 ID/OOD task (\text{0.86}).
    Tasks are scored directionlessly. \emph{At a
    strict $5\%$-FPR point mean ADV recall is \text{0.957} (median \text{0.964});
    Sec.~\ref{ssec:e2e} reports an end-to-end cascade analysis.}}}
  \label{fig:avg_auroc}
\end{figure}

\rev{Under the held-out 17-architecture protocol
(Table~\ref{tab:main_results}, Fig.~\ref{fig:avg_auroc}), Viyog leads on the
security task T2} (\textbf{0.980}) and
T3 (\textbf{0.844}) with $O(C)$ state. The five evaluated confidence/logit
scores (Energy, MSP, MaxLogit, Entropy, GEN) score at most $\text{0.647}$ on T2;
direction-consistent KL-Matching reaches $\text{0.659}$. A three-way supervised head reaches \textbf{0.798} balanced accuracy
(Fig.~\ref{fig:compl}). \rev{The T3 signal is above chance across the evaluated
17 usable architectures, four attacks, and three datasets
(Fig.~\ref{fig:t3_breakdown})}; GTSRB provides a second-ID-dataset check.


\begin{figure*}[tb]\centering
\includegraphics[width=0.85\textwidth]{figs/rebuttal/fig_complementarity.pdf}
\caption{\rev{\textbf{Combining a logit detector with Viyog yields three-way
ID/OOD/ADV separation: each covers the other's blind spot.} Logit detectors
(Energy) catch low-confidence OOD but are blind to high-confidence adversarials;
Viyog catches adversarials (dormant-band roughness) but is weak on smooth OOD.
\textbf{(a)}~Per-class recall (CIFAR-100, 17 archs): Energy alone $\text{0.445}$ on ADV,
Viyog alone $\text{0.257}$ on OOD; the \emph{Energy\,+\,Viyog} pair raises 3-way balanced
accuracy to $\mathbf{\text{0.758}}$ and the full panel to $\mathbf{\text{0.798}}$.
\textbf{(b)}~Gain transfers to GTSRB: full panel $\mathbf{\text{0.863}}$, logit$+$Viyog
$\text{0.768}$; bootstrap 95\% CIs are \emph{disjoint} (full-panel $[0.77,0.83]$ vs.\
Energy-only $[0.58,0.65]$). Viyog complements a logit OOD detector with no
retraining or extra forward pass, at $O(C)$ state.}}
\label{fig:compl}
\end{figure*}

\rev{The failure modes are complementary: logit scores collapse on T2} (blind to
high-confidence adversarials), while feature-distance detectors (Mahalanobis,
KNN) trail on T3, having pulled perturbations back toward the ID manifold at deep
layers~\cite{lee2018simple}. Among the compared single statistics, the
first-conv \emph{shape} statistic separates both (Sec.~\ref{ssec:rawbaseline}).

\begin{table}[tb]
\centering
\caption{\rev{\textbf{Held-out evaluation across the 17-architecture
CIFAR-100 panel} (of an original 20; MobileOne-S1, ResNet-18, and ResNet-34
excluded, Sec.~\ref{sec:setup}).
\emph{Setup:} T1/T2/T3 are defined in Sec.~\ref{sec:setup}.
\emph{Metric:} mean AUROC\,/\,FPR@95TPR (\%). $\times$~=~direction-inconsistent;
best AUROC per column in \textbf{bold};
\colorbox{viyoggreen!16}{shaded row}~=~primary detector; state memory in last column.
\emph{Highlights:} five of six logit baselines are direction-inconsistent on
T2; KL-Matching is consistent and reaches 0.659, vs.\ at most 0.647 for the
rest. Raw Viyog $L_\infty$ is itself direction-inconsistent, whereas
\textbf{Viyog} has the highest mean T2 (\textbf{0.980}) and T3 (\textbf{0.844})
among compared methods with $O(C)$ state.
\emph{Impact:} the evaluated signal is carried by first-conv
\emph{shape}, not magnitude or logits.}}
\label{tab:main_results}
\setlength{\tabcolsep}{3pt}
\renewcommand{\arraystretch}{1.08}
\footnotesize

\resizebox{\columnwidth}{!}{%
\revtab
\begin{tabular}{@{}lcccc@{}}
\toprule
 & \multicolumn{3}{c}{AUROC\,/\,FPR@95TPR (\%)} & \\
\cmidrule(lr){2-4}
Method & T1 ID/OOD & T2 ID/ADV & T3 OOD/ADV & State \\
\midrule
\multicolumn{5}{@{}l}{\textit{Logit baselines (\texttt{pytorch-ood})}}\\
Energy & 0.863 / 52.1 & 0.600\,$\times$ / 98.0 & 0.790 / 80.5 & 4\,B \\
MSP & 0.837 / 61.9 & 0.621\,$\times$ / 98.4 & 0.796 / 79.1 & 4\,B \\
MaxLogit & \textbf{0.864} / 52.5 & 0.599\,$\times$ / 98.6 & 0.794 / 79.5 & 4\,B \\
Entropy & 0.848 / 62.1 & 0.627\,$\times$ / 98.7 & 0.803 / 77.0 & 4\,B \\
GEN & 0.848 / 63.4 & 0.647\,$\times$ / 98.6 & 0.812 / 75.6 & 4\,B \\
KL-Match & 0.775 / 82.5 & 0.659 / 79.7 & 0.655 / 83.7 & 4\,B \\
\midrule
\multicolumn{5}{@{}l}{\textit{First-layer statistics (ours)}}\\
FirstConv-$L_\infty$ (baseline) & 0.562\,$\times$ / 98.1 & 0.515\,$\times$ / 87.6 & 0.577 / 75.0 & 8\,B \\
Viyog-HF (high-freq.\ dorm.) & 0.745 / 83.3 & 0.939 / 25.9 & 0.763 / 59.4 & $O(C)$ \\
\rowcolor{viyoggreen!16}\textbf{Viyog} (primary) & 0.663 / 96.5 & \textbf{0.980} / 4.9 & \textbf{0.844} / 37.4 & \textbf{$O(C)$} \\
\bottomrule
\end{tabular}}
\end{table}

\rev{A limited training-seed check (at most three seeds per architecture)
shows small T2 variability in this panel (Fig.~\ref{fig:seed_forest}); broader
training and data-split replication remains future work.}

\begin{figure}[tb]\centering
\includegraphics[width=\columnwidth]{figs/rebuttal/fig_seed_forest_new.pdf}
\caption{\rev{\textbf{Training-seed reproducibility across the full 20-architecture panel.}
Per-architecture mean~$\pm$~std (circles), individual seeds (faint crosses), panel
mean (vertical line), $\pm1$ overall std (band).
\textbf{(a)}~T2 seed std $\le0.004$ across all 20 backbones, an order of magnitude
below the cross-architecture spread (ViT-B the semantic-token outlier).
\textbf{(b)}~T3 similarly stable (std $\le0.06$).}}
\label{fig:seed_forest}
\end{figure}

\begin{figure*}[tb]\centering
\includegraphics[width=0.96\textwidth]{figs/rebuttal/fig_t3_breakdown.pdf}
\caption{\rev{\textbf{The OOD-vs-ADV (T3) signal is consistent across
architectures, attacks, and datasets.}
\textbf{(a)}~Per-architecture T3 AUROC on 17 usable CIFAR-100 backbones (mean
\text{0.844}, range \text{0.68--0.95}, every backbone above chance).
\textbf{(b)}~T3 across four attacks (FGSM, BIM, PGD, APGD-CE) and three datasets,
weakest case $\ge0.77$.
\textbf{(c)}~Viyog matches the strongest logit detector (GEN) within $\pm0.02$ at
$O(C)$ state and no extra forward pass, yet stays complementary to it
(Fig.~\ref{fig:compl}).}}
\label{fig:t3_breakdown}
\end{figure*}

\subsection{Near-OOD Difficulty and Shape-Depth Sweep}
\label{ssec:norm_sweep}
\rev{We break the T3 result by OOD difficulty and compare the shape read across
depth. In the evaluated splits, ID/OOD channels are generally smoother than
their adversarial counterparts (Fig.~\ref{fig:viyogd_sig}).

\emph{Harder near-OOD and texture splits.} Because near-distribution or
texture-biased OOD may themselves perturb early activations, we break the
17-architecture OOD-vs-ADV (T3) result by difficulty. Viyog reaches
far-OOD ($0.96$) and degrades gracefully but stays above chance on near-OOD
($0.74$ mean, up to $0.95$) and texture-OOD ($0.67$); near-OOD is indeed the
harder regime. A retrospective high-frequency diagnostic reaches $0.92$ on
texture-OOD where the total-variation read is weakest
(Fig.~\ref{fig:ood_difficulty}). Because the displayed panel selects the better
statistic per architecture using evaluation performance, it is an oracle upper
envelope, not an operational combination; a selection rule must be fixed on
disjoint calibration data. In the six-architecture, two-attack depth sweep, the shape read
peaks at depth~0 (mean T2~$0.965$, T3~$0.918$; the higher T3 than the
17-architecture main mean \text{0.844} reflects the easier subset) and decays with depth;
this experiment does not establish optimality for unseen architectures or
attacks (Fig.~\ref{fig:shape_depth}).}

\begin{figure}[tb]\centering
\includegraphics[width=\columnwidth]{figs/rebuttal/fig_shape_depth.pdf}
\caption{\rev{\textbf{The first conv is the strongest layer in the evaluated
depth sweep.} ID-vs-ADV (T2) AUROC of
the dormant-band shape statistic $V$ and of the raw $L_\infty$ statistic at every
depth, averaged over six architectures (ResNet-50, DenseNet-121, MobileNetV3,
ConvNeXtV2, Swin-T, ViT-B) and two attacks (PGD, APGD). The shape read (green)
peaks at the first conv (mean \textbf{0.965}) and decays with depth, whereas the
raw norm (red) is weakest there (\textbf{0.52}) and peaks only deep
(\textbf{0.80}, depth~6). The first-conv shape read beats the best raw-norm layer
at every evaluated depth.}}
\label{fig:shape_depth}
\end{figure}

\begin{figure}[tb]\centering
\includegraphics[width=\columnwidth]{figs/rebuttal/fig_ood_difficulty.pdf}
\caption{\rev{\textbf{Retrospective TV/HF diagnostic upper envelope.}
OOD-vs-ADV (T3) AUROC on the 17-architecture CIFAR-100 panel (mean~$\pm$~std),
split by OOD \emph{difficulty}: \emph{far} (a different dataset), \emph{near}
(a semantically similar shift), and \emph{texture} (texture-biased). The
primary dormant-band total-variation read (Viyog-D, green) is strongest on
far/near-OOD but weakest on texture; its high-frequency complement (HF, blue)
is the mirror image. The amber Panel takes the per-architecture better
test-set statistic and reaches \textbf{0.76}--\textbf{0.97}; it is an oracle
diagnostic, not a calibrated detector. It raises texture-OOD from
\textbf{0.67} to \textbf{0.92}, while
the raw $L_\infty$ magnitude norm (grey) stays near chance ($0.5$)
throughout: the separable signal is first-conv \emph{shape}, not magnitude.}}
\label{fig:ood_difficulty}
\end{figure}

\subsection{Computational Efficiency}
\rev{We next compare the accuracy and cost of the evaluated detectors.}

\begin{figure*}[tb]
  \centering
  \includegraphics[width=0.82\textwidth]{figs/rebuttal/fig_baselines_cost.pdf}
  \caption{\rev{\textbf{Standalone score-path runtime and stored state}
    (CIFAR-100; mean over ResNet-50, MobileNetV3, and DenseNet-121; H200 with
    CUDA synchronisation). \textbf{(a)}~These raw-input timings traverse
    different portions of the backbone: Viyog stops after the first layer,
    whereas logit and deep-feature scores require later layers. They therefore
    describe standalone score paths, not incremental second-stage latency or
    deployed speedup. The Viyog path is \textbf{0.05\,ms/img}.
    \textbf{(b)}~The plotted state comparison shows
    $0.23$\,KB for the plotted 64-channel profile against
    the $9$--$28$\,MB covariance/feature-bank of the distance detectors
    ($10^4$--$10^5\times$ smaller). Incremental overhead for a normal full
    inference is reported separately in Fig.~\ref{fig:onsilicon}.}}
  \label{fig:baselines_cost}
\end{figure*}

Figs.~\ref{fig:baselines_cost}--\ref{fig:onsilicon} and Table~\ref{tab:edge}
quantify the cost: Viyog uses $O(C)$ state (about 0.3\,KB at 64 channels) and
$2.28\%$ of CNN MACs with no extra forward pass. Its measured incremental
detector latency is 3.9\% of full inference on the H200; edge-energy values
below are roofline projections, not on-device measurements. The standalone
raw-input score-path timings in Fig.~\ref{fig:baselines_cost} are not deployed
speedup measurements.

\begin{table}[tb]\centering\footnotesize
\setlength{\tabcolsep}{4pt}
\caption{\rev{Second-stage detector cost
(architecture-derived). State is $O(C)$ for the first-layer detectors versus
$O(D^2)/O(ND)$ for feature-distance methods. The optional ``LDA head'' is a
tiny linear discriminant fit over the detector scores that turns the binary
detector into a 3-class ID/OOD/ADV classifier ($K{=}3$), adding ${<}0.1$\,KB and
no extra forward pass.}}
\label{tab:edge}
\revtab
\begin{tabular}{@{}lccc@{}}
\toprule
Detector & State & Extra fwd & Order \\
\midrule
Energy/MSP (logit) & 4\,B & 0 & $O(1)$\\
FirstConv-$L_\infty$ (baseline) & 8\,B & 0 & $O(1)$\\
\rowcolor{viyoggreen!16}\textbf{Viyog} (primary) & \textbf{$O(C)$} & 0 & $O(C)$\\
\quad$+$ LDA head (3-class, $K{=}3$) & $+{<}0.1$\,KB & 0 & $O(CK)$\\
Mahalanobis\,/\,ViM & 4.5--17\,MB & 1 & $O(D^2)$\\
KNN & 20--40\,MB & 1 & $O(ND)$\\
\bottomrule
\end{tabular}
\end{table}

\begin{figure}[tb]\centering
\includegraphics[width=\columnwidth]{figs/rebuttal/fig_onsilicon.pdf}
\caption{\rev{\textbf{Measured H200 detector cost and edge projection.}
\textbf{(a)}~Per-architecture cost of the
primary $V(x)$ detector as a fraction of full inference: the exact MAC fraction
(mean $2.28\%$) and the \emph{measured} batch-$1$ latency on an H200 (CUDA events;
mean $3.9\%$) show a single-digit-\% second-stage tax on this platform
silicon across CNNs, a ViT, and a ConvNeXt. \textbf{(b)}~Compute-bound roofline
projection of the detector's per-image energy on real edge accelerators (Jetson
Orin Nano, RPi5\,$+$\,Hailo-8L): $19$--$37\,\mu$J versus a $0.6$--$1.2$\,mJ full
inference. These compute-bound figures are optimistic lower bounds; memory,
runtime, and power effects require measurement on each target device.}}
\label{fig:onsilicon}
\end{figure}

Together the AUROC, AUTC, latency,
and memory results place Viyog favorably among the evaluated methods. We next
stress-test the detector against an
adaptive adversary.

% =============================================================
\section{Adaptive Attack Evaluation}
\label{sec:adaptive}

A rigorous adversarial evaluation must include an adaptive adversary who
knows the detection mechanism and specifically targets it.

\noindent\textbf{Threat model.} The attacker aims at \emph{both}
misclassification and detector evasion under $L_\infty$, $\varepsilon{=}8/255$.
Fig.~\ref{fig:adaptive_ladder} shows the full A0--A3 knowledge ladder; the
defender's hook, band $\mathcal{B}$, and threshold $\tau$ are fixed at deployment.

\subsection{Adaptive Attack Formulation}
A norm-aware attacker knowing that Viyog monitors first-layer $L_\infty$
norms would optimise perturbations to preserve these norms while inducing
misclassification. We maximise the following penalised PGD objective:
\begin{equation}
  \mathcal{L}_{\mathrm{adap}} =
  \mathcal{L}_{\mathrm{CE}}\bigl(f(x+\delta),y\bigr)
  - \lambda\,\bigl(\|f_0(x+\delta)\|_\infty - \|f_0(x)\|_\infty\bigr)^2,
  \label{eq:adaptive_loss}
\end{equation}
where $\lambda>0$ controls the penalty strength; we sweep $L_\infty$/$L_2$/combined
modes ($\varepsilon{=}8/255$, $\alpha{=}2/255$, 10 PGD steps). The reported
behavior is specific to these attacks, budgets, steps, and penalty sweeps.

\subsection{Results and Security Implications}

\rev{The ladder of Fig.~\ref{fig:adaptive_ladder} measures what each rung of
attacker knowledge costs, from black-box PGD (A0) through a both-aware adaptive
adversary (A3) to a stochastic-band defence (EOT).}

\begin{figure*}[tb]\centering
\includegraphics[width=0.96\textwidth]{figs/rebuttal/fig_adaptive_ladder.pdf}
\caption{\rev{\textbf{Defence-in-depth against an escalating adaptive adversary (mean over 20 architectures).}
Attacker knowledge grows left-to-right: A0~(PGD), A1~norm-preserving, A2~TV-aware
(white-box), A3~both-aware, EOT~stochastic defence.
\textbf{(a)}~A1 achieves \textbf{0\% evasion} empirically; A2 erodes $V$ to $0.64$ while
the complement holds at $0.74$; EOT restores $V$ to \textbf{0.82}; A3 suppresses
both only at an \textbf{8\% attack-success cost}.
\textbf{(b)}~EOT scatter: \textbf{18/20 architectures} are above the diagonal;
this supports stochastic band-randomisation in the tested setting but is not a
robustness certificate.}}
\label{fig:adaptive_ladder}
\end{figure*}

With the penalty active, $\|f_0(x_{\mathrm{adap}})\|_\infty\approx\|f_0(x)\|_\infty$,
keeping the raw score in-distribution. In the tested A1 configuration, the
dormant-band shape read has \textbf{0\%} measured evasion; this observation does
not follow from P1, because holding a norm fixed is not uniform rescaling. A3
reduces both evaluated signatures at a measured \textbf{8\%} attack-success
cost.

\emph{The result holds on a second dataset.} On GTSRB (ResNet-50,
ConvNeXtV2, ViT-B) at max penalty $\lambda$, the both-aware attacker keeps
$\ge\!80\%$ success yet the combined $\max(\textrm{dorm},\textrm{HF})$ detector
does not drop below AUROC $0.775$ in these runs: single-channel evasion is
backstopped by
the complement (ResNet-50 dorm $0.83\!\to\!0.51$ but HF stays $0.72$; ConvNeXtV2
dorm-suppression costs $80\%$ of success). The attacker thus either keeps the
detector $\ge\!0.775$ or surrenders most of its success, on both datasets.

\emph{Behavior of the evaluated attacks.} The tested standard
(A0) attacks~\cite{goodfellow2014explaining,madry2017towards,croce2020reliable,carlini2017towards,moosavi2016deepfool}
naturally suppress first-layer norms and are often detected; the bespoke
penalty of Eq.~\eqref{eq:adaptive_loss} targets that behavior directly. A
different white-box optimiser may find another trade-off, so these results do
not establish a universal cost of evasion~\cite{kim2020torchattacks}.

\emph{Failure-mode caveat.} Evading the global $L_\infty$ score may route a
sample toward ID or OOD. Abstention is safer only if the first-stage gate still
flags that sample; otherwise the adversarial can remain undetected.

\subsection{The Robustness--Evasion Trade-Off}
In the tested optimisation, norm preservation and misclassification compete,
producing an $8\%$ success cost for the both-aware attack. This is an empirical
frontier for the specified optimiser, not a fundamental lower bound.
A stronger held-out-basis (A4) attacker that re-optimises its perturbation
basis directly against the statistic shows a similar frontier
(Fig.~\ref{fig:adaptstrong}), with ${\sim}15\%$ of attack success surrendered
in those runs. We next test whether the observed phenomenon transfers beyond
the CNN architectures in the main panel.



% =============================================================
\section{Generalisation Beyond CNNs}
\label{sec:beyond_cnn}

The affine part of Proposition~\ref{prop:shape} applies to convolution and
patch-projection pre-activations, while its empirical premises still require
validation for each model. We evaluate three
architecturally distinct models fine-tuned on CIFAR-100 at $224\times224$
resolution from ImageNet-21k pretrained weights.

\textbf{ConvNeXtV2-Base}~\cite{woo2023convnext},
\textbf{ViT-B/16}~\cite{dosovitskiy2020image}, and
\textbf{Swin-T}~\cite{liu2021swin} span depthwise-convolutional, pure
self-attention, and hierarchical shifted-window designs; in every case the
first layer is an affine strided-convolution or patch-embedding map (ConvNeXtV2
$4{\times}4$ conv, ViT-B $16{\times}16$ and Swin-T $4{\times}4$ patch embeddings)
satisfying Eq.~(\ref{eq:affinebound}) exactly at the
\emph{pre-activation} stage. The implementation reads post-activation maps, for
which Proposition~\ref{prop:shape} does not provide a guarantee; the GELU
Taylor remainder is discussed in Appendix~\ref{app:affine}.

On all three, the dormant-band roughness $V(x)$ separates ADV from
ID/OOD as in the $32\times32$ CNN results, and the separation is
stable over the evaluated seeds and data resamples
(Fig.~\ref{fig:seed_forest}); at most three training seeds per architecture do
not establish seed invariance.



These results provide evidence that the signal transfers to several
convolutional and patch-projection architectures; they do not prove universal
architectural generality. We next evaluate a three-way classifier.

% =============================================================
\section{Downstream Utility: Three-Way Classification}
\label{sec:discussion}

Viyog turns a binary anomaly detector into a three-way ID/OOD/ADV classifier,
an operational distinction that diagnostic metrics (AUROC, FPR) do not capture:
without it a system must apply one response to all non-ID inputs, wasting
adversarial purification~\cite{nie2022diffusion,purify,silva2023diffdefense}
or feature-denoising~\cite{xie2019feature} on benign OOD, or leaving attacks
undetected. Beyond these binary routing cases, in
class-incremental learning~\cite{rebuffi2017icarl,cao2025class}
OOD inputs should feed back as new classes while ADV inputs must not contaminate
training.

\subsection{End-to-End Cascade and False-Positive Propagation}
\label{ssec:e2e}
\rev{We report a two-stage analysis (non-ID gate $\to$ Viyog) over the
17-architecture panel $\times$ 10 calibration/test split seeds at a matched
$5\%$ ID-FPR (Table~\ref{tab:e2e}, Fig.~\ref{fig:cascade_e2e}). Both stages'
thresholds are fit on each seed's calibration split only and evaluated once on
its disjoint test split (\texttt{heldout\_calibration\_v1},
Sec.~\ref{sec:setup}); this is the same protocol as the main results and
excludes MobileOne-S1, ResNet-18, and ResNet-34.}

\begin{figure}[tb]\centering
\includegraphics[width=0.95\columnwidth]{figs/rebuttal/fig_cascade_e2e.pdf}
\caption{\rev{Two-stage cascade over the 17-architecture panel at a
disjoint-calibration $5\%$ ID-FPR. \textbf{(a)}~Energy-only gate: mean
ADV recall $0.157$; two-signal gate ($\max(z_{\mathrm{energy}},|z_{\mathrm{dorm}}|)$, where $z_X$ denotes the normalised score for statistic $X$):
mean $0.834$. \textbf{(b)}~ID$\to$ADV mis-escalation is $\le\!0.0001$.}}
\label{fig:cascade_e2e}
\end{figure}

The evaluated Energy gate has low adversarial recall ($0.157$). Adding the
first-conv deviation to the gate
($\max(z_{\mathrm{energy}},|z_{\mathrm{dorm}}|)$, same $5\%$ ID-FPR) lifts it to
$\mathbf{0.834}$ (mean over all 17 architectures; the weakest is ViT-B
$0.29$ (its patch-embedding has no dormant low-activation band), then
ResNet-50 $0.70$, while OOD recall holds at $\approx0.42$)
with mean ID$\to$ADV mis-escalation of $0.0001$ (exactly $0$ on $15/17$
architectures).

\begin{table}[tb]\centering\footnotesize
\setlength{\tabcolsep}{4pt}
\caption{\rev{Two-stage recall at a disjoint-calibration $5\%$ ID
false-positive rate (mean over 17 archs $\times$ 10 split seeds).}}
\label{tab:e2e}
\revtab
\begin{tabular}{@{}lccc@{}}
\toprule
Stage-1 gate & ADV recall & OOD recall & ID$\to$ADV\\
\midrule
Energy only & 0.157 & 0.343 & 0.002\\
Energy $\oplus$ dorm-$z$ (2-signal) & \textbf{0.834} & 0.424 & \textbf{0.000}\\
\bottomrule
\end{tabular}
\end{table}


\begin{figure}[tb]\centering
\includegraphics[width=0.84\columnwidth]{figs/rebuttal/fig_viyogd_signature.pdf}
\caption{\rev{\textbf{The dormant-band signature.}
Per-image dormant-band TV profile for ID (green), OOD (blue), and adversarial
(red) inputs on ResNet-50 (CIFAR-100, five OOD sets, PGD attack). ID and OOD
share a low, smooth profile; adversarials inject high-frequency structure,
separating at AUROC~$0.91$.}}
\label{fig:viyogd_sig}
\end{figure}

\begin{figure}[tb]\centering
\includegraphics[width=0.88\columnwidth]{figs/rebuttal/fig_adaptive_strong.pdf}
\caption{\rev{\textbf{Strong adaptive frontier: held-out basis (6 archs).}
An A4 attacker that re-optimises its perturbation basis directly.
Dorm/HF mean \textbf{0.86}/\textbf{0.79} (min $0.54$/$0.53$); only the
both-aware objective suppresses both to $\approx$\textbf{0.62} at a
\textbf{15\%} attack-success cost in this evaluated configuration.}}
\label{fig:adaptstrong}
\end{figure}

\begin{figure}[tb]\centering
\includegraphics[width=0.82\columnwidth]{figs/rebuttal/fig_nonvision.pdf}
\caption{\rev{\textbf{Synthetic 1D sanity check.}
ID-vs-ADV AUROC of $V(x)$ (green) vs.\ raw $L_\infty$ (red) at the first
\texttt{Conv1d} layer for FGSM/PGD at two budgets ($n{=}1{,}500$, 3 seeds).
Shape reaches $0.95$--$1.00$ while the raw norm stays near chance; this single
synthetic task is preliminary evidence, not a general non-vision validation.}}
\label{fig:nonvision}
\end{figure}
\subsection{Limitations and Future Work}
\label{ssec:limitations}
\rev{Table~\ref{tab:limits} states the principal scientific limitations. In
particular, the affine proposition is conditional and the adaptive evaluation
is not a certificate.}

\begin{table}[tb]\centering\footnotesize
\caption{\rev{\textbf{Current limitations and required follow-up.}}}
\label{tab:limits}
\setlength{\tabcolsep}{3pt}
\revtab
\begin{tabular}{@{}p{1.5cm}p{5.6cm}@{}}
\toprule
Limitation & Current boundary and required follow-up \\
\midrule
(i) Theory & Conditional affine pre-activation bound; nonlinear hooks are
empirical, with no worst-case guarantee. \\
(ii) Threat & Four digital $L_\infty$ attacks at $8/255$ and finite adaptive
sweeps; test other norms, physical attacks, and unseen objectives. \\
(iii) Accuracy & T3 FPR@95 is \text{37.4}\%; cascade OOD recall is 0.424 and ViT ADV
recall is 0.29; these are not stand-alone safety evidence. \\
(iv) Generality & Three image datasets, $\le3$ seeds/model, one synthetic 1D
task. Excludes MobileOne-S1/ResNet-18/ResNet-34 (zeroed features,
Sec.~\ref{sec:setup}); $\dagger$-marked studies (Table~\ref{tab:coverage})
predate this and still report 20. \\
(v) Hardware & H200 latency is measured; edge energy is projected and
requires measurement on target devices. \\
\bottomrule
\end{tabular}
\end{table}

% =============================================================
\section{Conclusion}
\label{sec:conclusion}

We proposed Viyog, a gradient-free method that uses first-layer dormant-band
roughness to separate OOD from ADV inputs. Under the evaluated digital
$L_\infty$ threat model, Viyog reaches T2 AUROC $\mathbf{\text{0.980}}$ across the
17-architecture held-out panel. Proposition~\ref{prop:shape} gives a sufficient
pre-activation separation condition; it is not a guarantee for arbitrary
attacks, natural shifts, or nonlinear hooks. In the adaptive suite, A1 has
$0\%$ measured evasion, while TV-aware attacks reduce performance.

\rev{Combined with logit scores, Viyog reaches three-way balanced accuracy
\textbf{0.798} on CIFAR-100 and \textbf{0.863} on GTSRB, adding no learned
parameters and $O(C)$ state at 2.28\% CNN MAC / 3.9\% H200 latency overhead.
High T3 FPR, weak cases, limited adaptive coverage, and projected edge energy
prevent stand-alone deployment claims without the follow-up in
Table~\ref{tab:limits}.}

\appendices
\section{General-Affine Bound, GELU Remainder, and Proof of Proposition~1}
\label{app:affine}
The first-layer perturbation bound (Eq.~\eqref{eq:affinebound}) uses only that the first layer is
\emph{affine}. For any first-layer Jacobian $W$ (dense layer, strided
convolution, or ViT/Swin patch embedding) and perturbation $\|\delta\|_\infty
\le\varepsilon$,
\begin{equation}
  \|W\delta\|_\infty \le \|W\|_\infty\,\varepsilon,\qquad
  \|W\|_\infty=\max_i\textstyle\sum_j|W_{ij}|,
  \label{eq:affinebound}
\end{equation}
the induced $\infty$-operator norm (maximum absolute row sum). Convolution is
the special case in which $W=M_m$ is Toeplitz and the maximum row sum equals
$\|W_m\|_1$, recovering Eq.~\eqref{eq:affinebound}; the Toeplitz structure is thus
\emph{sufficient but not necessary}, so the mechanism extends to patch-embedding
first layers that have no shift-invariance.

If the first activation is a smooth $g$ (e.g.\ GELU), the linearisation error
in each pre-activation coordinate is bounded by the Taylor remainder
$\tfrac12\sup_z|g''(z)|\,(\|W\|_\infty\varepsilon)^2$. For
$\mathrm{GELU}(z)=z\,\Phi(z)$, $g''(z)=\varphi(z)\,(2-z^2)$ has interior extrema
at $z=0$ (value $2\varphi(0)$) and $z=\pm2$ (value $-2\varphi(2)\approx-0.108$);
since $0.798>0.108$ the supremum is at $z=0$,
\begin{equation}
  \sup_z|\mathrm{GELU}''(z)| = 2\varphi(0) = \sqrt{2/\pi} \approx 0.798,
\end{equation}
so the per-coordinate error is at most $0.399\,(\|W\|_\infty\varepsilon)^2$.
Normalising this quadratic bound by the perturbation envelope
$\|W\|_\infty\varepsilon$ gives
$0.399\,\|W\|_\infty\varepsilon$, which is $O(1)$--$O(10)$ for the measured
first-layer weights (e.g.\ ${\approx}6.8$ for ResNet-50 at
$\varepsilon{=}8/255$). This is not a relative-error guarantee because the
actual linear response may be smaller than its envelope. The calculation
therefore does not certify a post-GELU version of
Proposition~\ref{prop:shape}; it only quantifies a worst-case local remainder.

\subsubsection*{Full Proof of Proposition~\ref{prop:shape}}
The affine residue $e_c=w_c\!\ast\delta$ established above is the object on which
the roughness statistic acts; we now prove its three properties (P1--P3) in turn.

\noindent\textbf{(P1) Uniform rescaling.} Absolute 1-homogeneity gives
\[
\widetilde{\mathrm{TV}}(\lambda_c a_c)
=\frac{\lambda_c\mathrm{TV}(a_c)}{\lambda_c\nu_c+\epsilon}
=\frac{\mathrm{TV}(a_c)}{\nu_c+\epsilon/\lambda_c}.
\]
This is exact invariance only when $\epsilon=0$; it makes no claim about
perturbations that preserve a norm while changing spatial structure.

\textbf{(P2) Two-sided bound.} Affineness gives
$e_c=w_c\!\ast\delta$ exactly. TV subadditivity and reverse triangle give
\[
[\rho_c r_c-\beta_c]_+\nu_c
\le \mathrm{TV}(a_c+e_c)
\le(\beta_c+\rho_c r_c)\nu_c.
\]
The spatial $L_1$ triangle and reverse-triangle inequalities give
\[
\nu_c|1-r_c|\le\overline{|a_c+e_c|}
\le\nu_c(1+r_c).
\]
For the lower ratio bound, divide the lower numerator by the upper denominator;
for the upper ratio bound, divide the upper numerator by the lower denominator.
Adding $\epsilon=\eta_c\nu_c$ and cancelling $\nu_c$ yields P2.

\textbf{(P3) Conditional separation.} A natural comparison satisfying A1 has
normalised TV at most $\beta_c$. Hence P2 gives
\[
\widetilde{\mathrm{TV}}(a_c')-\widetilde{\mathrm{TV}}(a_c)
\ge
\frac{\rho_c r_c-\beta_c(2+r_c+\eta_c)}
     {1+r_c+\eta_c}.
\]
The stated condition makes each margin positive, and averaging over the fixed
band $\mathcal B$ yields the claim. The result is conditional on A1--A3 and
does not extend automatically through ReLU or GELU.





\section{Non-Vision Sanity Check (1D Signals)}
\label{app:nonvision}
As a preliminary check beyond images, we evaluate a 1D-CNN trained on
band-limited multi-sinusoid signals. It detects adversarials at mean T2
AUROC~$\mathbf{0.96}$ (vs.\ $0.71$ for the raw norm, 3 seeds;
Fig.~\ref{fig:nonvision}). This synthetic experiment is evidence for one 1D
setting, not a claim of modality-agnostic performance.

\section*{Artifact Availability}
The evaluated MIT-licensed source, tests, paper source, and reproduction entry
point are archived as release \texttt{v0.1.3-codes26-ae-6} at
Zenodo~\cite{viyogartifact2026} and mirrored by the exact
\href{https://github.com/amanyagami/viyog/releases/tag/v0.1.3-codes26-ae-6}{GitHub release}.
The detector is distributed as
\href{https://pypi.org/project/viyog/0.1.3/}{\texttt{viyog==0.1.3}}, and the
checkpoints are pinned at
\href{https://huggingface.co/amanyagami/viyog-weights/tree/b96263ae28a8585acaad6a002a0298b6b6c2d735}{Hugging Face}.
The canonical entry point is \texttt{reproduce\_t1.sh}; the precomputed
adversarial dataset and live Space linked from the archive README are optional
and are not required by that script.

\section*{Acknowledgment}
This research was supported by: the Global Learning \& Academic research
institution for Master's\,\textperiodcentered\,PhD students, and Postdocs
(G-LAMP) Program of the National Research Foundation of Korea (NRF), funded
by the Ministry of Education (No. RS-2025-25442252); an NRF grant funded by
the Korea government (MSIT) (RS-2025-00561169); an NRF grant funded by
the Korean government (MSIT) (No. RS-2022-00165225, RS-2026-25494138); the Institute of
Information \& Communications Technology Planning \& Evaluation (IITP)
under the Leading Generative AI Human Resources Development program
(IITP-2026-RS-2026-25546026), funded by the Korea government (MSIT); the
Ministry of Science and ICT (MSIT), Korea, under the National Program in
Medical AI Semiconductor (No. 2024-0-0097), supervised by IITP in 2026; an
IITP grant funded by the Korean government (MSIT) (No. RS-2025-02304331,
Digital Columbus Project); National Science Foundation grant ACED 2436016;
the ANCHOR Program through the Gangwon ANCHOR Center, funded by the
Ministry of Education (MOE) and Gangwon State (G.S.), Republic of Korea
(No. 2026-RISE-10-006); and the BK21 FOUR project (AI-driven Convergence
Software Education Research Program), funded by the Ministry of Education,
School of Computer Science and Engineering, Kyungpook National University,
Korea (No. 41202420214871).

OpenAI Codex was used to assist with language editing, LaTeX/typesetting
checks, and code/test auditing during camera-ready preparation. The authors
reviewed and take responsibility for all technical claims, equations,
experiments, and final text.

\bibliographystyle{IEEEtran}
\bibliography{refs}

\end{document}
